% Part: normal-modal-logic
% Chapter: frame-correspondence
% Section: first-order-definability

\documentclass[../../../include/open-logic-section]{subfiles}

\begin{document}

\olfileid{nml}{frd}{fol}

\olsection{First-order Definability}

我们已经看到，框架可达关系的许多性质都可以用模态公式来定义。例如，框架的对称性可以用公式~\Ax{B}，即 $p \lif \Box\Diamond p$ 来定义。迄今为止我们遇到的条件，都可以在包含单个二元谓词符号的语言中用一阶公式表达。例如，对称性由
$\lforall[x][\lforall[y][(\Atom{Q}{x,y} \lif \Atom{Q}{y, x})]]$ 定义，其含义是：一个论域为 $\Domain{M} = W$ 且 $\Assign{Q}{M} = R$ 的一阶结构~$\Struct{M}$ 满足上述公式，当且仅当 $R$ 是对称的。这提示了如下定义：

\begin{defn}
  一个框架类~$\mClass{F}$ 是\emph{一阶可定义的}，如果存在带有一个二元谓词符号~$Q$ 的一阶语言中的语句~$!A$，使得 $\mModel{F} =
  \tuple{W, R} \in \mClass{F}$ 当且仅当在论域为 $\Domain{M} = W$ 且 $\Assign{Q}{M} = R$ 的一阶结构~$\Struct{M}$ 中，$\Sat{M}{!A}$ 成立。
\end{defn}

事实上，迄今为止所讨论的性质及定义这些性质的模态公式都属于特例。并非每个公式都能定义一阶可定义的框架类，也并非每个一阶可定义的框架类都能由某个模态公式定义。

第一种主张的反例由 L\"ob 公式给出：\begin{equation}
\Box(\Box p \lif p) \lif \Box p. \tag{\Ax{W}}
\end{equation}
\Ax{W} 定义了传递且逆良基的框架类。一个关系是良基的，如果不存在无限序列 $w_1$, $w_2$, \dots{} 使得 $Rw_2w_1$, $Rw_3w_2$, \dots。例如，$\Nat$ 上的 $<$ 关系是良基的，而 $\Int$ 上的 $<$ 关系则不是。一个关系是逆良基的，当且仅当其逆关系是良基的。因此，逆良基关系是指不存在无限序列 $w_1$, $w_2$, \dots{} 使得 $Rw_1w_2$, $Rw_2w_3$, \dots 的关系。

然而，不存在定义传递且逆良基关系的一阶公式。假设存在一阶公式~$!F$ 使得 $\Sat{M}{!F}$ 当且仅当 $R = \Assign{Q}{M}$ 是传递且逆良基的。令 $!A_n$ 为公式
\[
(\Atom{Q}{a_1,a_2} \land \dots \land
    \Atom{Q}{a_{n-1},a_{n}})
\]
现在考虑公式集
\[
\Gamma = \{!F, !A_1, !A_2, \dots\}.
\]
$\Gamma$ 的每个有限子集都是可满足的：设 $k$ 是使得 $!A_k$ 属于该子集的最大下标，构造 $\Domain{M_k} = \{1, \dots, k\}$，
$\Assign{a_i}{M_k} = i$，以及 $\Assign{Q}{M_k} = <$。由于 $\{1, \dots, k\}$ 上的 $<$ 是传递且逆良基的，所以 $\Sat{M_k}{!F}$。根据构造，$\Sat{M_k}{!A_i}$ 对所有 $i \le k$ 成立。根据一阶逻辑的紧致性定理，$\Gamma$ 在某个结构~$\Struct{M}$ 中可满足。根据假设，既然 $\Sat{M}{!F}$，关系 $\Assign{Q}{M}$ 就是逆良基的。但显然，$\Assign{a_1}{M}$, $\Assign{a_2}{M}$, \dots{} 将构成一个逆良基性所禁止的那种无限序列。

第二种主张的反例是全域性（universality）性质：对所有 $u$ 和 $v$，都有 $Ruv$。全域框架可由公式 $\lforall[x][\lforall[y][\Atom{Q}{x,y}]]$ 一阶定义。然而，不存在在所有且仅在全域框架中有效的模态公式。这是一个本身也很有趣的结果的推论：在全域框架中有效的公式，与在自反、对称且传递的框架中有效的公式完全相同。存在非全域的自反、对称且传递的框架，因此，每个在所有全域框架中有效的公式，也在某些非全域框架中有效。

\end{document}